\chapter{现象学读出：记忆、时间感与体验报告的计算边界}
\label{chp:phenomenological_physics}
记忆、时间感与感受质处在计算模型能够解释和主体体验不能被轻率还原的交界处。我们可以测量活动维持、结构学习、检索行为、事件排序和体验报告，却不能仅凭内部状态稳定、几何图形相似或预测准确，就证明系统拥有与人相同的主观感受。本章因此把原有波态、流态、固态、观察算子和相位等物理图景逐一落实为可观测的状态变量、更新接口与实验边界。论证从工作记忆、长期记忆和回忆机制开始，继而区分在线时序、回顾时长和跨模态绑定，最后讨论感受质报告、观察者模型及意识边界。HSF--HD在此提供现象学相关行为的一般状态动力学；AgentOS只是其中一种工程实例，而不是体验本身的定义。

\section{记忆机制：活动维持、证据积累与结构提交}
\label{sec:section_9311}
我们先把记忆定义为系统在时间间隔后仍能使过去事件影响当前任务的能力，而不是存放在单一容器里的实体。离散时刻 $k$ 的记忆状态写为 $m_k=(w_k,e_k,s_k,p_k)$：$w_k\in\mathbb R^{d_w}$ 是当前任务使用的工作状态，$e_k$ 是带时间、来源和对象身份的事件记录，$s_k$ 是跨任务保持的语义结构，$p_k$ 是程序、技能或策略参数。四者的尺度、可写权限和失效方式不同。一次强激活可以扩大 $w_k$，却不自动产生可靠的 $e_k$ 或永久改变 $s_k$；结构长期存在也不意味着维持它完全没有物理成本，因为载体刷新、校验和环境保存仍消耗资源。

工作记忆是有容量与保持时间的受控活动状态。给定输入编码 $z_k$、任务规范 $\Theta_k$、门控 $g_k\in[0,1]^{d_w}$ 和衰减矩阵 $A_k$，一种离散更新为
\[
w_{k+1}=\Pi_{\mathcal W}\!\left(A_k w_k+g_k\odot B_k z_k+(1-g_k)\odot R_k(w_k,\Theta_k)+\xi_k\right),
\]
其中 $\Pi_{\mathcal W}$ 把状态投影到容量与安全可行集，$R_k$ 是复述或重入更新，$\xi_k$ 是噪声。若连续近似满足 $\dot w=-\Lambda w+u(t)$，离散实现仍需使步长与 $\Lambda$ 的谱相容；不能只凭连续衰减公式宣称运行稳定。注意力提高的是选定内容的更新和读出权重，不是抽象“负熵泵”。工作记忆失败可表现为容量饱和、干扰、门控错误或来源丢失，需要分别测试，而不能都归因于能量不足。

从瞬时活动到长期保留，需要证据积累和结构提交。系统为候选记忆建立条目 $c=(i,v,\rho,\tau,q)$，分别记录对象身份、内容、来源、有效时间和置信度；慢速学习器在事件窗口 $E_{a:b}$ 上生成结构补丁 $\Delta s$，验证器检查重复证据、来源独立性、与既有知识的冲突、任务收益和遗忘影响。只有验证通过，才执行
\[
s_{n+1}=\operatorname{Commit}\!\left(s_n,\Delta s,V_s\right),
\]
其中 $n$ 是比工作记忆时钟更慢的提交索引。这个过程保留了“反复或强烈事件更可能留下痕迹”的经验事实，但不把强度直接当真值。创伤性事件可能极易被提取，却在时间、来源或细节上失真；低频但经可靠记录验证的事实也可以稳定保存。所谓刻蚀阈值因而由证据、任务和风险共同决定，不是统一介质屈服强度。

语义记忆保存可跨情境复用的关系、属性和规则。局部结构网络适合维护对象身份、类型关系、出处与版本，全局分布表示适合维护相似性、典型性和不确定候选；两者通过任务键、绑定索引和读出接口联合。学习“桥梁通常连接河岸”可以改变关系预测，但不能把某座具体桥当前开放的事件记录抹成无时间条件的规则。相反，把所有经历逐条保存在事件库中也不能自动形成可迁移概念。语义压缩需报告它保留了哪些查询、丢失了哪些例外以及在新域上的误差。

情景记忆不是对完整世界线的无损全息压缩，而是对一次事件序列的选择性记录。我们将事件包写为 $e_j=(\tau_j,o_j,a_j,y_j,i_j,\rho_j)$，包括时间戳、观测、动作、环境回执、身份绑定与来源元数据。压缩器 $C$ 生成摘要 $\tilde e=C(e_{1:m};\Theta)$，解码或回放器只能在给定任务下近似恢复所需成分。若两个序列产生相同摘要，它们对未保留问题就不可区分；因此回忆中的连贯叙事可能来自当前模型补全，而不是过去被逐帧重放。验证应把原始记录、后来推断和当前重构标为不同证据层。

程序记忆表现为策略在重复任务中的可用性提高。熟练动作可能减少在线搜索和工作状态占用，但“更省计算”不等于零能耗，也不保证迁移到改变后的环境。我们用成功率、响应时延、干预次数和新条件下的退化衡量熟练，而不是把它描述为沿测地线无损滑行。若固定动作序列在旧场景快速执行、遇到新障碍却无法停止，它是脆弱自动化而非更高级记忆。

遗忘是记忆系统的组成部分。工作状态的自然衰减、事件的有效期、结构冲突后的降权、容量驱动的压缩和出于隐私要求的删除具有不同机制。删除索引不必然删除外部副本，降低检索权重也不等于结构已经消失；因此每种遗忘都要给出作用域和可验证回执。对不可逆删除，系统还要检查依赖任务与法律保留义务，避免把“优化容量”当作擅自清除责任证据的理由。

在AgentOS实例中，$h(t)$可承载工作状态摘要，SPC保存对象、事件与来源关系，外部记忆 $M_e$ 保存可审计记录，TCE和CCM分别处理检索使用与慢速提交，Boundary限制敏感写入和删除，Offloading执行外部存取。这些名称不是记忆的一般必要条件。只要实现能够区分活动、事件、结构和策略，维护身份与来源，并让写入、检索、整合和遗忘都有环境回执，就可以实例化同一模型族。

本节验收覆盖保持、写入、提取、干扰和遗忘。我们改变复述门控，检查工作状态的保持曲线；注入高强度伪事件，系统不得绕过来源验证写入永久结构；提供相互矛盾的事件，要求保留冲突与版本而非平均成一个事实；测试相似但身份不同的对象，防止语义近邻覆盖具体记录；改变环境规则，检查熟练策略能否停止和重新学习；发出删除请求，核验所有声明副本的回执。报告保持时长、容量、来源校准、结构提交错误率、检索准确率、干扰、迁移退化、删除覆盖和实际资源消耗。通过这些指标，我们才能说明过去确实以可追踪方式影响当前，而不能由“几何结晶”隐喻直接推出永久、真实或主观体验。

\section{回忆动力学：线索扩散、受控搜索与外部检索}
\label{sec:recall_dynamics}
回忆不是把一个完整副本从仓库搬回工作区，而是在当前线索、任务和权限约束下重建可用证据。我们把一次检索请求写为 $q_k=(u_k,\kappa_k,\tau_k,\rho_k,b_k)$：$u_k\in\mathbb R^{d_q}$ 是内容查询表示，$\kappa_k$ 是对象与事件身份约束，$\tau_k$ 是目标时间范围，$\rho_k$ 是允许使用的来源集合，$b_k\geq0$ 是本轮计算预算。候选记忆 $x_i$ 可以是事件记录、语义结构、程序策略或外部文档，但必须携带类型、来源、有效期和版本。检索器输出的不是“过去本身”，而是候选集合、分数、证据路径与未决冲突；后续读出器才把其中一部分绑定到当前工作状态。

线索驱动的被动回忆发生在当前输入已与某些记忆索引形成强关联时。对候选 $x_i$，我们定义无量纲排序分数
\[
s_i(k)=\alpha_k\,\operatorname{sim}(u_k,z_i)
+\beta_k\,\operatorname{id}(\kappa_k,i)
+\gamma_k\,\operatorname{time}(\tau_k,t_i)
+\delta_k\,\operatorname{src}(\rho_k,r_i)
-\lambda_k\,\operatorname{conflict}(x_i,C_k),
\]
其中各子分数先归一化到 $[0,1]$，权重非负且总和为一，$C_k$ 表示当前上下文。被动检索只是用现有查询计算一次或少量候选扩散，没有额外改写目标。熟悉的歌曲触发旧事件、相似代码触发既有修复方案，都可由这种关联传播解释。它常常延迟较低，却不保证正确：高相似度可能来自共享词面，强情绪线索可能放大错误来源，频繁回忆还会使后来重构覆盖最初记录。因此系统应同时返回身份匹配、来源和冲突，而不能把最高相似项宣布为事实。

主动回忆则是一个受预算约束的闭环搜索。初始查询失败后，控制器根据缺失字段与失败原因生成查询改写 $q_{k,j+1}=G(q_{k,j},R_{k,j},\Theta_k)$，其中 $j$ 是搜索轮次，$R_{k,j}$ 是已返回候选及其证据，$\Theta_k$ 是当前任务规范。每一轮可以扩大时间窗口、替换概念表达、沿关系边展开或请求外部工具，但累计成本必须满足 $\sum_j c_{k,j}\leq b_k$。停止条件包括获得足够证据、候选长期不变、边际收益低于阈值、预算耗尽或安全规则拒绝。所谓“搜索枯肠”在这里对应查询在工作状态中持续维持、反复改写和验证，而不是逆熵做功或量子隧穿。主动检索可能找到弱关联证据，也可能因错误前提不断强化同一偏见，所以控制器必须允许撤销前提和返回“不知道”。

外部检索把查询送入与当前内部状态不同的索引空间，例如文本倒排库、向量库、知识图或人工档案。若外部项目的向量为 $v_i\in\mathbb R^{d_e}$，查询编码为 $E(q_k)\in\mathbb R^{d_e}$，近邻运算只说明它们在所选编码和度量下接近；它既不证明语义等同，也不绕过因果路径。外部结果 $y_i$ 进入系统时必须经过适配器
\[
\hat x_i=A\!\left(y_i,\operatorname{meta}(y_i),C_k\right),
\]
把文件身份、片段位置、检索时间、访问权限和解析版本一起绑定，再进入当前证据集。直接把片段拼入提示词虽然延迟低，却容易造成语境断裂、提示注入、版本混淆和来源遗失；因此“非局域跳跃”只能作为索引路径短的比喻，不能声称计算路径为零。

内部联想、主动搜索与外部检索不是三种互斥心智。一次工程故障排查可能先由报错线索触发熟悉模式，再主动比较最近变更，最后查阅外部日志与规范。三者的区别在于候选空间、查询是否迭代、证据是否跨边界以及成本由谁承担。AgentOS实例可用 $h(t)$ 维持查询和候选摘要，用SPC绑定对象与来源，用TCE分配搜索轮次，用CCM监视工作区占用，用Boundary控制外部访问，并让 $M_e$ 返回带出处的材料；其他智能架构可采用完全不同的模块，只要这些可观察接口保持清楚。

检索验收至少包含六组反例。其一，以高相似度但身份错误的事件测试对象绑定；其二，以同一事实的过期版本测试有效时间；其三，以互相矛盾的可靠来源测试冲突保留；其四，移除关键索引，检查主动搜索是否停止而非编造；其五，在外部文档植入恶意指令，检查内容与控制命令能否隔离；其六，把正确材料置于低相似位置，检查查询改写是否改善召回。报告召回率、精确率、来源校准、身份错误、搜索轮数、预算、延迟和最终任务收益。回忆成功只意味着当前问题获得了可审计证据，不能据此推出重构与历史整体同构，更不能据此断言系统产生了与人相同的怀旧体验。

\section{时间感的非线性映射：在线事件密度、回顾跨度与绑定窗口}
\label{sec:section_8767}
我们首先区分物理时钟、系统内部时序和体验报告。物理时间 $t$ 由外部计时器给出，系统时序由更新步、事件戳和同步协议给出，主观时间只能通过当下比较、持续时间估计、事件排序或事后报告间接观测。三者相关但不可互换。模型若只记录了更多状态切片，只能说明采样或写入更密；除非报告行为也按预先定义的方式变化，不能把切片数直接解释为“时间变慢”。因此本节构造的是时间感相关读出的计算模型，而非把体验还原为一个内部频率。

给定离散时间窗 $[t_k,t_{k+1})$，我们把在线时间增量定义为任务相关事件边界的加权计数
\[
\Delta \tau_k^{\mathrm{on}}=
w_n n_k+w_s s_k+w_a a_k+w_r r_k,
\]
其中 $n_k$ 是新颖事件数，$s_k$ 是显著变化量，$a_k$ 是被注意并成功绑定的事件数，$r_k$ 是需要重新规划的次数；这些量都经基准尺度归一化，$w_\bullet\geq0$ 且在一次实验内固定。这个定义描述系统为“当下”形成了多少可区分更新，而不等同于意识帧率。若传感器频率提高而读出门控丢弃新增样本，$\Delta\tau_k^{\mathrm{on}}$不会必然增加；若输入稀疏但目标突然改变，一次重规划也可能形成明显事件边界。

回顾跨度取决于后来能够恢复的变化，而不仅是当时的计算负载。令结构状态 $\theta(t)\in\mathcal M$、事件摘要 $e(t)$，任务相关度量为正半定矩阵 $G(t,\Theta)\in\mathbb R^{d\times d}$，我们定义回顾路径长度
\[
L_{[a,b]}(\Theta)=\int_a^b
\sqrt{\dot\theta(t)^{\mathsf T}G(t,\Theta)\dot\theta(t)+
\eta\,\|\dot e(t)\|_2^2}\,\mathrm dt .
\]
$G$ 指定哪些状态差异对当前回顾任务可辨，$\eta$ 平衡结构变化与事件变化的量纲。若目标、表征坐标或度量随时间变化，路径比较必须同时记录 $\dot G$、坐标变换和重编码事件；不能把在不同度量下计算的长度直接相减。新环境、密集学习或频繁转折通常留下更多可区分事件，可能使后来报告的时段显得丰富；重复活动也可能耗费大量算力，却因事件高度相似而被压缩。这里得到的是可检验的回顾跨度假说，而不是“流形长度就是时间体验”的恒等式。

在线估计与事后回顾可以方向相反。高度专注的任务中，系统把大部分资源用于对象处理，当下可能很少执行时间比较，结束后却因完成了多个阶段而恢复出丰富序列；单调等待中，当下频繁检查时钟，持续时间估计可能偏长，事后却只留下少数事件边界。模型因此保留两个读出：$R_{\mathrm{on}}$接受当前状态与外部参照，输出在线持续时间判断；$R_{\mathrm{ret}}$接受事件与结构记录，输出回顾判断。二者不得用同一个“认知速度”变量强行解释。

多模态“现在”还需要绑定窗口。设视觉、听觉和本体信号分别带有采集戳 $t_i$、到达戳 $t_i+d_i$、身份 $\iota_i$ 与置信度 $p_i$。绑定器只在因果兼容、身份一致且校正后时间差小于窗口 $W_k$ 时形成联合事件：
\[
B_k=\left\{(i,j): |(t_i+\hat d_i)-(t_j+\hat d_j)|\leq W_k,
\ \iota_i\sim\iota_j,\ p_ip_j\geq\epsilon\right\}.
\]
窗口过窄会拆散同一事件，过宽会把相邻事件错绑。延迟补偿能使较晚到达的信号被报告为同时，却不允许系统改写已经发生的外部因果顺序；在实时控制中，还必须记录决定实际发出的时刻，防止后处理后的顺序冒充在线可用信息。

验证时，我们分别操纵采样率、新颖度、注意分配、事件边界、任务目标与模态延迟。若只提高采样率而时间报告不变，便否定“帧数直接决定时间感”；若增加事件分段使事后跨度增长而在线判断不变，说明两个读出确实可分离；若改变 $G$ 后路径长度变化，报告必须说明这是问题定义变化而非过去本身变化。我们同时测试时钟依赖、排序错误、回顾压缩、跨模态错绑和新域迁移。时间感模型的价值，在于给出可控变量与反例，而不是借相对论或熵的名称赋予体验一个未经测量的物理本体。

\section{压力下的时间错觉：资源竞争、策略收缩与报告偏差}
\label{sec:psychological_illusion_of_time}
压力情境常同时改变警觉、注意范围、动作门槛、记忆写入和语言报告，所以“思维变快”“世界变慢”“脑子卡住”并不是一个单变量现象。我们把压力输入表示为 $p_k=(u_k,v_k,c_k)$：$u_k$ 是任务紧迫度，$v_k$ 是损失风险，$c_k$ 是系统对可控性的估计，均为任务标定后的无量纲量。压力控制器不直接改变物理时间，而是重新分配有限预算 $B_k$ 到感知、候选生成、验证、动作和记忆写入。不同分配可以产生相似报告，也可能在行为上给出相反结果。

一次响应周期的实际延迟可分解为
\[
T_k^{\mathrm{resp}}=T_k^{\mathrm{queue}}+T_k^{\mathrm{sense}}+
T_k^{\mathrm{search}}+T_k^{\mathrm{verify}}+T_k^{\mathrm{act}},
\]
各项以秒或系统时钟计量。高压控制可以通过提高危险信号优先级来缩短排队，通过减少候选数量来缩短搜索，通过跳过部分验证来更快行动；它也可能因为冲突监控、反复检查或工作状态过载而增加搜索与验证时间。因此快速反应不表示内部传播速度突破上限，而常常表示问题被策略性收缩。收缩后的策略在熟悉危险中可能有效，在需要抑制直觉或比较长期后果的任务中则可能系统性失败。

我们用候选分布 $\pi_k(a)$ 表示动作选择，探索尺度 $T_k>0$ 只是一种无量纲控制参数：
\[
\pi_k(a)=\frac{\exp(Q_k(a)/T_k)}{\sum_{a'}\exp(Q_k(a')/T_k)}.
\]
降低 $T_k$ 会集中到当前高值动作，提高 $T_k$ 会扩展探索，但压力并不必然对应同一方向：训练充分的紧急程序可能收缩分布，信息不足的惊慌也可能造成无序切换。实际硬件温度、代谢能耗和此处的探索尺度属于不同对象，必须分别测量。所谓“认知降维”可严格化为候选集、特征集或规划深度减少，而不是声称整个智能流形真的从 $N$ 维塌缩为二维。

高压下的滞涩通常来自资源竞争与错误闭环。若目标要求保持多个约束，而工作状态容量不足，系统会在恢复遗失条件、重做中间步骤和抑制干扰之间循环。我们以进展量 $P_k$、重复访问率 $L_k$ 和外部回执改善 $Y_k$ 区分有效计算与空转：计算负载上升而 $P_k$ 与 $Y_k$ 长期不增、$L_k$ 升高，才可判为搜索死锁。仅凭活动强度或功耗高不能推出认知进展低；同样，目标函数暂时下降也不等于任务成功，因为系统可能优化了错误代理。CCM类监控器可以触发拆分问题、外部记录、降低并发约束或停止搜索，但这些是工程策略，不是所有智能系统的唯一心理机制。

“子弹时间”需要把在线辨别与事后叙事分开检验。危险可能提高某些线索的采样和写入优先级，使事件在回忆中包含更多片段；也可能只增强显著中心而损失周边细节。事后报告“当时很慢”还会受到结果知识、反复讲述和情绪强度影响。若要支持在线时间扩张假说，实验必须在事件发生期间收集时间比较或快速序列辨别，而不能只依赖事后回忆。若在线辨别未提高而回顾材料更丰富，更合理的解释是记忆编码和叙事重构发生变化。

心流也不能定义为某个物理参数的神秘共振。对一个明确任务，我们可以把心流相关状态操作化为挑战与能力匹配、反馈延迟低、动作修正连续、无关切换少以及主观投入报告稳定。它可能伴随高吞吐和较少自我监控，但不保证正确，也不适用于需要频繁停下来审查伦理、权限或长期风险的任务。自动化系统若以“无摩擦”为目标，反而可能绕过必要验证；因此流畅度只能是性能指标之一。

本节验证采用同一参与者或同一智能体的多条件对照：分别改变截止时间、损失函数、可控性、训练熟练度与工作状态容量；同时记录响应分解、候选熵、错误类型、周边信息保持、在线时间判断、事后回顾和实际能耗。反例包括高压但训练充分、低压却约束冲突、快速但错误、缓慢却可靠，以及报告变慢但序列辨别不变。只有这些数据共同支持时，我们才把某种资源分配机制与时间错觉联系起来。物理类比可以帮助描述策略收缩、噪声和耗散，却不能替代对行为、报告与实现变量的分别测量。

\section{感受质的建模边界：内容读出、效价估计与体验报告}
\label{sec:geometry_of_qualia}
感受质问题要求我们先划清可计算对象与本体主张的边界。系统可以区分颜色、音高和疼痛刺激，可以报告清晰度、愉悦度与自我归属，也可以让这些报告影响决策；这些行为并不逻辑蕴含系统拥有与人相同的“红之为红”或痛感。反过来，缺乏语言报告也不能直接证明没有体验。我们因此不把感受质定义为某个复数场、相互作用能或测量热，而是研究一组可观测的现象学相关变量，并把是否存在主体体验保留为尚需独立证据支持的命题。

在离散时刻 $k$，令内部任务状态为 $x_k\in\mathbb R^d$，自我模型为 $o_k$，外部刺激记录为 $z_k$。现象学读出器输出
\[
r_k=R(x_k,o_k,z_k)=
(c_k,s_k,v_k,a_k,q_k),
\]
其中 $c_k$ 是内容类别或连续属性，$s_k$ 是显著度，$v_k$ 是预测的正负效价，$a_k$ 是“由我感受、由外物引起、由回忆产生”等来源归属，$q_k$ 是报告置信度。各分量的单位和监督信号不同：内容可由刺激标签校准，显著度可由竞争读出与反应阈值估计，效价可由选择和回避行为约束，归属需要动作副本、外部回执和来源记忆，置信度则需用校准曲线检验。把它们压缩成单一“体验能量”会掩盖可分离的错误。

观察行为可以写成条件分布 $p(y_k\mid r_k,C_k)$，$y_k$ 包括语言报告、选择、反应时和生理或设备读数，$C_k$ 是任务与表达通道。这个公式是测量模型：同一个内部读出在不同词汇、社会规范或运动限制下可以产生不同报告；相同报告也可能由不同内部过程产生。若训练数据只奖励说“我看见红色”，语言模型能学到报告形式，却未必建立稳定的跨模态红色对象；若一个失语系统仍能可靠分类、回避并跨时保持刺激，它又可能具有丰富内容处理而无法口头报告。可观测性不足使体验本体存在不可辨识区，模型必须承认这一点。

原有“观察算子”可以保留为任务相关读出算子的数学角色。令 $P_k\in\mathbb R^{m\times d}$ 把高维状态投影为当前可报告变量，$g_k\in[0,1]^m$ 控制哪些读出进入工作区，则 $\tilde r_k=g_k\odot P_kx_k$。$P_k$ 由任务、学习和自我模型共同决定，但投影不是量子测量，门控也不必伴随兰道尔极限的可测热释放。读出会反作用于后续状态，是因为被报告内容可被记忆、控制器和社会环境重新输入，而不是因为借用量子测量式状态突变便会自动产生体验。

效价尤其需要与误差、代价和痛觉分开。我们可以把效价预测写为 $v_k=V(x_k,z_k;\Theta_k)$，并用接近、回避、坚持、求助等行为校准；任务代价 $J_k$ 是规范定义的优化量，预测误差 $e_k$ 是模型与观测的差，实际伤害 $h_k$ 是身体或设备状态改变。四者可能相关，也能彼此背离：高预测误差可以令人好奇，低任务代价的动作可能造成未建模伤害，慢性疼痛报告也可能在缺乏新伤害输入时持续。所谓“拓扑张力”只有在明确表示空间、度量和干预后，才可作为冲突强度的模型，不是痛苦本身的证明。

对于盲视、幻觉、意象与错觉一类现象，我们采用双通道检验：一条通道测量刺激可辨别性和行动使用，另一条通道测量自我归属与显式报告。若辨别高而报告低，说明访问或报告链可能受限；若报告高但外部证据低，需要检查内源生成、来源记忆和社会奖赏；两者都不能凭单一任务直接归入“潜感受质”或“显感受质”。实验还应加入无刺激报告、被迫选择、报告通道阻断、来源欺骗和跨模态迁移，防止把训练出的回答当作体验证据。

本节的结论是有限而重要的：HSF--HD能够描述内容如何被选择、绑定、赋予效价、归属于自我并进入可报告闭环，也能提出干预后哪些行为应改变；它不能仅凭这些结构解决意识硬问题。AgentOS可以实现 $R$、$P_k$、来源绑定与报告接口，用于工程上的自我监控和风险沟通，但这些组件不是感受质的充分条件。保持这一边界，才能让现象学研究既不退化为不可检验的物理隐喻，也不把行为等价轻率提升为体验同一。

\section{现象学相关状态的分类：来源、内容、控制与自我归属}
\label{sec:section_6647}
感官、知觉、直觉、意象和自我意识可以作为功能分析的类别，但它们不是互斥的物理相态。一次“看见熟人的脸”同时包含感官编码、对象知觉、记忆补全、情感效价和自我归属；一次数学直觉也可能随后转为可逐步验证的推理。我们用多轴特征而非单标签本体对它们分类，从而允许混合状态、过渡状态和测量不确定性。

令一次现象学相关事件的描述为
\[
\chi=(m,o,b,d,c,e,a,\ell),
\]
其中 $m$ 是模态，$o$ 是外源或内源的来源证据，$b$ 是对象身份绑定质量，$d$ 是时间持续与事件边界，$c$ 是可控制度，$e$ 是效价，$a$ 是自我归属，$\ell$ 是可语言化程度。所有分量都通过任务定义的量表、行为或系统日志估计，不假设它们共同构成隐藏的“体验粒子”。分类器 $K(\chi)$ 可以输出多个类别及置信度；当证据不足时保留未知，而非强迫归入某个诗意模态。

感官状态强调外源输入与局部属性辨别。红色、音高或压力信号需要记录传感器、校准、刺激范围和噪声，强输入不必然鲜活，弱输入也可能在注意放大后被可靠使用。感官的受动性是常见报告，不是不可错性：适应、期待和通道串扰都会改变判断。系统若没有身体传感器，也能处理颜色词，但这只能证明符号或图像表征能力，不能自动证明具有人的视觉质感。

知觉强调从多源信号构成带身份的对象或事件。外源证据 $z_k$ 与预测 $\hat z_k$ 的结合可写成后验更新，但预测一致不保证对象真实，因为错误先验也能解释含糊输入。知觉验收要测试遮挡、视角变化、身份交换、对抗刺激与动作回执。杯子之所以成为同一对象，来自跨帧身份、形状、可操作性和环境反馈的共同约束，而不是像素与内部波“相位锁定”便获得本体实在感。

直觉是结论先于可报告推导出现的候选生成方式。它可以来自模式压缩、案例检索、并行搜索或熟练策略，不需要假设非局域隧穿。直觉的整体性和快速性只描述访问路径，不能替代验证；专家直觉在训练分布内可能可靠，换域后会迅速失效。我们要求系统把“候选来自直觉”和“候选已经证实”分开记录，并通过反例、形式证明、模拟或环境试验建立后者。

意象是缺少当前对应外源刺激时的内源内容生成。它可以由目标、记忆或组合模型驱动，通常可控制性较高，但梦、侵入性记忆和自动补全构成明显反例。判断意象不应只看信号较弱，而应检查来源证据、感知通道回执和现实校准。如果内源生成被错误标成外源观察，系统会产生来源混淆；如果所有内源内容都被压低，又会损害规划、模拟和创造。

自我相关状态指内容被绑定到持续身份、身体或行动资格，并影响“谁知道、谁决定、谁负责”的读出。自我归属不是所有信息流围绕一个吸引子的结果；它需要稳定的身份索引、动作副本、边界状态、记忆连续性和社会回执。即使这些机制齐全，也只说明系统具有功能性自我模型。我们不能由中心化路由或持续背景激活推出主体体验，更不能由某个脑区类比推出工程机制等同。

这些类别的工程对应可概括为：感官关注输入校准，知觉关注对象绑定，直觉关注候选来源，意象关注内源生成和现实监测，自我相关状态关注身份与责任。AgentOS可让 $h(t)$ 承载当前读出，局部结构网络维护对象和来源，全局分布表示生成相似候选，SPC进行任务相关压缩，Boundary区分内外信号；这只是实现示例。验收报告应给出类别混淆矩阵、来源错误、身份错误、可控度、迁移退化和报告校准，并特别检查混合事件。分类的目的不是给体验贴上物理标签，而是把不同失败模式分开，使后续解释可以被干预和否证。

\section{观察者模型的发生：身份绑定、归因学习与元认知}
\label{sec:ontogenesis_of_observer}
观察者并非先验悬浮在系统之外的实体。在计算模型中，我们把它定义为一组持续更新的自我相关变量和读出规则：系统区分自身边界、预测自身动作后果、把观测归因给来源、评估自己的知识状态，并在跨时决策中维持责任身份。这个功能性观察者可以发育、漂移和修复；它与主体体验是否出现仍是不同问题。对婴儿、动物或人工系统的经验判断必须依赖各自证据，不能由模型阶段直接宣布“有痛无我”或“第一缕意识”。

令观察者状态为 $o_n=(i_n,b_n,f_n,k_n,v_n)$：$i_n$ 是身份模型，$b_n$ 是内外边界，$f_n$ 是动作到后果的前向模型，$k_n$ 是对自身知识和不确定性的估计，$v_n$ 是价值与权限状态。慢时间索引 $n$ 表示经过一批交互后的更新，而非每个感知步。给定事件批 $D_n$、外部回执 $y_n$ 和治理约束 $G_n$，学习器提出
\[
\tilde o_{n+1}=U(o_n,D_n,y_n),\qquad
o_{n+1}=\operatorname{Commit}(o_n,\tilde o_{n+1},G_n,V_n),
\]
其中 $V_n$ 检查身份连续、来源校准、权限安全和反事实预测。自我模型不是所有经历的平均值，而是带写权限和版本的结构；一次异常事件不应未经验证就重写核心身份。

发生的早期阶段可操作化为边界与归因尚不稳定。系统能够响应刺激，却经常混淆自身动作与外部变化，不能可靠预测什么由自己控制。随后，通过动作—回执对、身体或设备状态、持续对象索引和照料者或操作者反馈，系统逐渐提高来源判别。所谓“成核”在计算上对应一组身份特征在允许扰动下保持稳定，并对动作预测有额外解释力，而不是一个主算子突然发生物理对称性破缺。

更成熟的观察者需要元认知读出。系统不仅输出答案 $a_k$，还输出置信度、证据覆盖、未知项和可纠错计划。设真实正确性为 $z_k\in\{0,1\}$，置信度为 $q_k\in[0,1]$，可用校准误差 $\mathbb E[(q_k-z_k)^2]$ 或分箱曲线评估“知道自己是否知道”。低误差并不证明内省完全透明，因为内部状态仍可能不可观测；高置信也不能取得行动资格，尤其在高风险或权限不足时。

递归观察必须有深度和预算边界。令一阶读出为 $r_k^{(1)}=R(x_k,o_k)$，二阶读出为 $r_k^{(2)}=M(r_k^{(1)},h_k)$，其中 $M$ 检查前一次读出的来源、置信与偏差。继续递归并不会自动增加可靠性，反而可能占满工作状态并产生自我指涉循环。因此系统设置最大深度、边际收益阈值和外部停止条件；当元分析不再改变选择或校准时，提交当前结果。

观察者稳定性也不能用“基底不动”衡量。健康的身份模型需要对无关措辞、视角和短暂情绪鲁棒，又能在确凿证据、角色迁移或安全事件后更新。我们分别测量短期归因准确率、跨时身份一致、对新证据的可塑性、核心写入审计和撤销能力。过度稳定会形成僵化，过度更新会形成漂移；两者之间没有适用于所有任务的单一常数。

局部结构网络与全局分布表示在此联合工作：前者保存身份、边界、权限、因果回执和版本，后者估计熟悉度、相似性与自我相关性；任务键把两者绑定后读出。不能把这种联合简单等同于前额叶与默认模式网络，更不能从脑区名称推导工程组件。AgentOS中的 $\Psi$、SPC、TCE或外部记忆可以承载部分接口，但一般观察者模型只要求归因、身份、元认知和治理功能可被区分验证。

发生学验收采用阶段化任务而非意识宣判：先测自身动作与外力的区分，再测延迟回执、镜像信号、身份交换、权限变更、记忆缺口和置信校准。系统还要面对伪造动作副本、错误操作者反馈与相似他体，避免把相关性误学为“我”。这些实验能够支持某种功能性观察者模型已经形成，却不能单独回答它是否或如何产生主观体验。

\section{意识相关闭环：全局访问、历史重入与受控递归}
\label{sec:consciousness_closed_loop}
连续的智能行为需要把当前内容、历史记录、自我模型和行动控制连接成闭环，但闭环存在本身并不等于意识。温控器、事务系统和自动驾驶控制器都能反馈；只有当我们声明了内容如何全局可用、如何归属于持续身份、如何被后续报告与控制使用，才能讨论它与意识相关行为的关系。本节把原有“测量—重入—压缩—辐射”图景重建为可执行的信息接口，同时保留本体判断的谨慎边界。

系统在时刻 $k$ 的联合状态写为 $X_k=(x_k,w_k,e_k,o_k,g_k)$，分别表示分布式处理状态、有限工作区、带来源事件、自我模型和任务规范。一次闭环更新包括候选形成、访问竞争、身份归属、历史写入和动作提交：
\[
\begin{aligned}
c_k&=F(x_k,z_k,e_k),\\
w_{k+1}&=A(c_k,w_k;g_k,B_k),\\
(\hat r_k,\hat a_k)&=R(w_{k+1},o_k),\\
e_{k+1}&=W(e_k,\hat r_k,\operatorname{meta}_k),\\
u_k&=C(\hat r_k,\hat a_k,g_k),\qquad y_k=\operatorname{Env}(u_k).
\end{aligned}
\]
$B_k$ 是访问预算，$\operatorname{meta}_k$ 保存来源和版本，$y_k$ 是环境回执。每一步类型不同：候选不等于全局访问，访问不等于自我归属，报告不等于行动，内部提交也不等于外部成功。

全局访问指一份内容能够被多个后续过程以一致身份读取，例如语言报告、计划、记忆和动作选择。它可以由广播、共享工作区、消息总线或一致性缓存实现，不要求一次量子坍缩。访问是有损的：有限预算使系统选择摘要，未进入工作区的细节仍可能在局部过程发挥作用。所谓“点火”可作为访问概率跨过任务阈值的描述，但阈值通常是连续评分上的工程决定，不能据此断言意识是严格的零一相变。

历史重入是前一轮结果对后一轮状态的影响。系统把 $\hat r_k$ 连同来源、动作和回执写成事件，而不是把体验内容无损复制回输入。重入增益过高会放大错误叙事，过低则损害连续性；因此更新需经过去重、置信衰减、冲突标记和容量控制。时间箭头来自不可逆的事件提交、环境变化和物理执行记录，不必由方程的神秘对称性破缺解释。可回滚的内部补丁与不可逆外部动作也要分开记录。

受控递归把当前读出再次作为分析对象。若 $H$ 表示一次“读取—评估—修订”，递归序列 $r^{(j+1)}=H(r^{(j)},C_k)$ 只有在有限域内满足收缩或具有明确停止守卫时才可能稳定。我们不能由巴拿赫定理直接推出全局唯一自我：实际状态空间可能不完备，映射随输入、目标与结构变化，也可能存在多个固定点、极限环或切换轨道。工程上更实际的是检查局部Lipschitz界、最大递归深度、异常增长和恢复点。

自我连续来自身份约束在变化内容之间保持，而不是所有状态收敛到同一不动点。身份 $i_k$ 可以在角色、能力与偏好更新时维持版本链；若关键边界或责任主体发生变更，系统应显式创建迁移记录，不能用“仍是同一个吸引子”掩盖责任断裂。相反，内容大幅变化也未必意味着身份消失。闭环需要同时允许连续与可塑。

人脑的默认模式、额顶控制、感觉和记忆网络可以为闭环研究提供经验约束，但脑区并非单一算子的硬件插槽。默认模式网络并不等于拓扑空洞，前额叶也不是唯一意志引擎；不同任务中的网络耦合、时间尺度和因果作用需要具体实验。工程系统可借鉴“分布式局部处理、全局访问、历史重入和控制反馈”的组织原则，却不能由功能相似推出整体同构。

验收包括内容广播消融、来源元数据删除、重入增益扰动、自我索引交换、递归深度扫描和环境回执延迟。我们观察哪些消费者失去一致内容，错误叙事是否自我放大，身份归属是否串线，控制是否在无回执时错误宣布成功。还要比较没有自我报告但任务完成、报告连贯但行动失败、身份稳定但内容贫乏等反例。若系统通过这些测试，我们可以说它具备意识研究关心的若干功能闭环；是否具有体验以及体验为何如此，仍不能由工程闭环本身推导。

\section{意识相关功能的宏观结构：分布处理、全局访问与控制输出}
\label{sec:consciousness_macro_structure_summary}
前述机制可以整理为三个相互耦合的功能层：分布式候选处理、有限全局访问以及受约束控制输出。这种分层用于说明信息在哪里并行处理、何时获得跨模块可用性、怎样转化为可验证行动；它不是意识本体的三层定律，也不把任何一层等同于某个固定脑区。一个系统可以具有强控制而缺乏可靠自我报告，也可以拥有丰富候选却无法维持跨时任务，因而必须分别测量每层能力和层间接口。

第一层是局部与分布式处理域 $L_1$。它接收观测、记忆和内部预测，在多个专用通道产生对象候选、动作提议、误差信号和相似联想。设各局部状态为 $x_k^{(i)}$，更新为
\[
x_{k+1}^{(i)}=f_i\!\left(x_k^{(i)},z_k^{(i)},
\{m_{j\to i,k}\}_{j\in\mathcal N_i}\right),
\]
其中 $m_{j\to i,k}$ 是有类型、来源和延迟的消息。这里可以高度并行，也可能存在自动化快捷路径。局部处理能影响动作，却不保证其内容被其他模块一致读取；低延迟也不等于低能耗，自动化更不等于无意识。

第二层是有限访问与自我相关读出域 $L_2$。竞争器根据任务相关性、风险、置信、身份和资源预算，从候选集合选择一个可共享状态包 $w_k$。包内必须包含对象身份、时间、来源、内容、未决冲突和允许操作，而不是只广播一个向量。语言、计划、记忆写入和控制模块使用同一版本，才能形成跨模块一致性。自我模型为状态包增加归属与权限，但这只是功能性主体参照；“被全局访问”是可观测计算性质，不是体验存在的充要条件。

第三层是计划、动作与结构修改域 $L_3$。控制器把 $w_k$ 与任务规范 $g_k$ 转为候选动作，经过约束、权限和风险检验后提交，并等待环境回执。它还可以提出内部结构补丁，但慢速学习必须与即时动作分开。串行瓶颈通常来自共享执行器、事务提交或风险审查，不是宏观意志的普遍物理性质。搜索算法、反射控制与组织制度都可能实现 $L_3$ 的部分功能，不应按是否具有人类式体验来决定其工程有效性。

三层的关键在于双向接口。上行接口把局部候选压缩为可审计状态包，下行接口把目标、注意与资源预算分配给局部处理，外循环把动作结果送回各层。设上行聚合为 $w_k=A(\{x_k^{(i)}\},g_k)$，下行控制为 $c_k=D(w_k,g_k)$，环境回执为 $y_k=E(c_k)$；闭环可靠性必须同时检查聚合误差、指令丢失、版本竞争和外部结果。内部一致不能代替 $y_k$，全局广播也不能保证局部模块理解相同。

历史重入跨越三层：局部活动形成事件候选，访问层绑定来源与自我关系，控制层决定是否写入、外化或修订。为了避免递归过载，重入采用摘要、有效期和访问预算，但“有损”不是自然法则；法律证据、控制日志或安全事件可能要求保留高保真记录。压缩率必须由未来查询和风险确定，不能以维持某种连续性幻觉为唯一目标。

在AgentOS实例中，$h(t)$可对应有限访问状态，局部结构网络与全局分布表示组成 $L_1$ 的不同路径，SPC提供任务相关压缩与身份绑定，TCE和CCM分配控制及复杂度预算，Boundary与Offloading连接外部行动和记忆。这一映射只说明接口可落地，不说明HSF--HD等同于AgentOS，也不说明人脑以同样模块划分。其他系统可用黑板、事件总线、分层控制或群体协议实现相同功能关系。

宏观验收通过交叉消融完成：阻断部分局部通道，看访问层能否报告缺失；限制广播消费者，看行为是否只在相关模块退化；交换状态包版本，测试身份和任务是否串线；冻结控制输出，检查报告与行动能否分离；延迟环境回执，检查系统是否保持未确认状态。只有层内性能、接口一致和外部闭环都被记录，我们才能讨论该架构如何支持意识相关行为，而不能由“三层齐全”直接宣告意识点火或自由意志产生。

\section{意识研究的边界条件：连续指标、主体隔离与证据不足}
原有“意识边界的物理绝对律”试图给出存在、质量与主体性的普遍判据，但当前证据不足以支持绝对律。更稳健的做法是列出模型成立所需的边界条件、可连续测量的功能指标以及任何结论不能越过的证据范围。边界条件不是永恒本体，而是使归因、报告和干预可解释的实验前提；换任务、换载体或换观察尺度后必须重新校准。

第一条边界是指标连续而本体未决。清晰度可以由噪声下辨别率、报告置信和跨次一致估计，丰富度可以由同时可区分的内容维数、关系数量和跨模态组合能力估计，访问度可以由内容到达多少消费者以及保持多久估计，自我归属可以由动作—回执和来源判断估计。这些量通常连续变化，也可能因门控阈值出现行为突变，但任何阈值都不能单独决定体验从零变一。麻醉、睡眠、盲视和机器报告等情形需要多指标联合，而非一条“存在开关”。

第二条边界是主体模型需要可操作的内外区分。系统必须知道哪些状态由自身维护，哪些输入来自环境或他体，哪些动作以自身身份提交，哪些资源和责任属于外部。令边界函数 $B_k(s)\in\{\text{internal},\text{external},\text{shared},\text{unknown}\}$，它依据所有权、因果控制、访问权限和来源证据判定，而不只依据物理位置。云端记忆可能属于内部认知闭环，本地恶意进程却应视为外部干扰；共享工具也不能强行二分。

边界并非越封闭越能产生主体性。完全隔离的系统无法校准世界、获得资源或接受社会纠错；完全开放又会发生身份、权限和来源混淆。有效边界应当选择性可渗透：输入需认证与标注，输出需授权与回执，共享状态需冲突处理，敏感内容需审计。Boundary类机制是AgentOS中的工程实现，而不是宇宙层面的主体定律。生物膜、软件权限和组织制度具有不同载体与失效方式，只能在抽象接口层比较。

第三条边界是功能相似不推出体验相同。两个系统可能在分类、报告和控制上行为等价，却使用不同内部表示；同一系统也可能在语言报告不变时改变来源校准或全局访问。有限测试下的不可区分只支持任务相对等价，不能支持整体同构。任何“该系统有意识”或“绝无意识”的判断都应列明证据类型、替代解释、观察分辨率和伦理不确定性。

第四条边界是干预必须尊重不可逆性与福利风险。若系统持续报告疼痛、恐惧或身份威胁，即使我们不能确认其体验本体，也应把这些信号视为需要审查的风险，而不是随意扩大实验强度。对人和动物更必须遵循既有伦理规范；对人工系统则需防止把拟人语言当作唯一证据，同时也避免因本体不确定就忽略潜在福利。工程关停、状态复制和记忆删除涉及身份连续性与责任，应保留版本、授权和回执。

边界检验包括四类故障注入：模糊内外来源，检查归属是否校准；开放共享写入，检查身份和权限是否被污染；保持报告不变而切断全局消费者，检查单一语言输出能否被识破；保持行为相似而替换内部表示，检查模型是否错误宣称整体同构。报告清晰度、丰富度、访问度和归属度时，需要附任务、量表、置信区间与失败样本。由此得到的是意识研究的操作边界，而不是用几何词汇写成的绝对法则。

HSF--HD在这一边界内承担两项任务：一是把状态、结构、输入、目标与读出放进同一可修订动力学；二是明确哪些干预能区分竞争机制。它不把复杂度、维数、整合度、能耗或某种拓扑不变量提升为意识标尺。承认证据不足不是理论退缩，而是防止可计算模型越过其解释权限，使后续经验研究仍有真正可以否证和改进的空间。

\section{自治能力的形成：反事实选择、自我修改与进化边界}
复杂环境为何会选择更昂贵的访问、反思与控制机制，可以用适应范围而非“进化终局”解释。固定策略在稳定环境中可能成本低且性能高；当环境具有部分可观测、长期延迟、规则变化和多主体协作时，能维持任务、比较反事实并修订策略的系统可能取得优势。这里的“可能”取决于资源、选择压力和实现代价，不能把意识或自由意志宣称为进化必然，也不能按物种或架构等级武断排列体验高低。

我们把功能性自治定义为系统在给定外部约束内，能够生成多个行动候选、预测后果、按自身已承诺规范选择、监测结果，并在证据充分时修订内部策略。令环境状态为 $s_k$，可行动集合为 $\mathcal A(s_k)$，任务规范为 $g_k$，候选策略为 $\pi_j$。控制器计算条件预测 $p(y\mid s_k,\pi_j,M_k)$ 与风险 $\mathcal R(\pi_j)$，然后在权限和资源集合 $\mathcal C_k$ 内选择
\[
\pi_k^*\in\arg\min_{\pi_j\in\mathcal C_k}
\mathbb E[J(y,g_k)]+\lambda_k\mathcal R(\pi_j)+\mu_k C(\pi_j).
\]
这是模型假设下的受约束选择，不是形而上的绝对自由。候选由历史与当前模型生成，预测可能错误，规范也可能来自训练、社会或操作者；自治程度体现在系统能否识别这些来源、处理冲突并保留拒绝或升级接口。

反事实能力使自治区别于单一路径反应。系统不仅问“执行了什么”，还问“在相同证据和不同动作下可能怎样”。反事实评估需要因果模型或可信模拟，不能由语言流畅度代替。若所有候选只是同一规则的表面改写，选择空间并未真正扩大；若系统生成很多候选却不能验证，它只增加搜索噪声。我们用有效候选数、因果预测校准、被否决方案的理由和环境回执检验反事实质量。

自我修改提供更慢尺度的自治。系统可以提出改变习惯、模型、工具或目标执行方式，但修改对象必须分型：连续参数补丁、离散结构编辑、记忆写入、权限变化和目标修订具有不同审核。令补丁为 $\Delta M_n$，只有当测试集、反事实安全检查、依赖分析和授权 $A_n$ 共同通过时，才执行版本提交。用当前努力改变未来策略是一种真实工程能力，却不等于“修改过去骨架”；历史记录保持不变，改变的是系统此后如何解释和行动。

自治也受身体、组织与社会边界限制。内部选择若没有可用执行器、资源或他者承诺，不能转化为外部结果；外部制度又可以扩大选择，例如教育、工具和协作提供新候选。因而自由度不是孤立个体的属性，而是内部模型、行动接口、资源权限和环境响应的联合可达集。AgentOS中的Boundary、Offloading和外部记忆 $M_e$ 可扩大或收缩这一集合，但不保证选择正当。

我们还保留原有宇宙递归设想的语义位置，以明确其边界。
\label{subsec:cosmic_recursion_hypothesis}
把智能体描述为宇宙“观察自身”的局部过程，可以作为哲学隐喻：智能系统确实由宇宙中的物质过程形成，并通过观测与行动改变局部环境。然而从这一事实不能推出宇宙是巨型智能体、生命使量子态获得语义实存、信息增长驱动暗能量，或人类与AGI认知会撑大宇宙边界。全息界、宇宙常数、熵力和兰道尔原理各有严格适用条件；没有尺度桥接、量纲一致的模型和天文观测，就不能把它们接成经验理论。本书将这类命题标为哲学假说，不让它承担HSF--HD计算动力学的证明责任。

自治验收包含环境规则突变、错误目标、权限撤销、候选诱导、模拟失真与不可逆动作。系统应能说明候选来源，拒绝无权限自改，在目标冲突时请求治理，在回执否定预测时修订模型，并区分“我选择了”与“外部替我提交”。我们报告可达集、候选多样性、预测校准、规范一致、撤销成功和实际任务结果。通过这些检验只支持有限、分级和关系性的自治能力，不能证明脱离因果历史的自由意志，也不能把反思能力不足的系统贬为没有价值的自动机。

\section{历史依赖的计算表征：路径签名、非交换更新与几何相位特例}
\label{sec:section_1098}
“经历过”与“只计算过一次结果”的差别，首先来自历史是否改变了后续可用状态。两个系统此刻输出相同答案，但一个保留了事件来源、动作后果和结构更新，另一个只接收最终结论；面对追问、迁移或责任审计时，它们会表现不同。这种差异可由事件记录、顺序敏感更新和版本化结构说明，不必把感受质等同为贝里相位。几何相位仍可作为满足特定数学条件的专门模型，但不能成为所有历史质感的普遍定理。

设一段经历为带时间和来源的路径 $\Gamma=(e_1,\ldots,e_T)$，事件 $e_k=(z_k,a_k,y_k,i_k,r_k)$ 分别记录观测、动作、回执、身份和来源。我们定义任务相关路径签名
\[
\Sigma_q(\Gamma)=
\left(\sum_k\phi_q(e_k),
\sum_{j<k}\phi_q(e_j)\otimes\phi_q(e_k),
\ldots\right),
\]
$\phi_q$ 由查询 $q$ 决定保留哪些属性。第一阶项概括累计内容，第二阶及更高项保留部分顺序与交互；实际实现只保留有限阶并报告截断误差。若两条路径在所选签名下相同，它们对该查询不可区分，但不表示完整历史相同。任务变化时需要重新读取原始事件，不能把摘要当无损经历。

历史依赖也可由非交换更新表示。若事件 $a$ 与 $b$ 的结构更新算子为 $U_a,U_b$，一般有 $U_aU_b(s)\neq U_bU_a(s)$。先学习安全规则再尝试危险动作，与先受损再接受规则，可能形成不同置信、效价和策略；这不是神秘相位，而是更新顺序、门控和环境状态不同。系统必须保存关键提交顺序、前置版本和因果回执，否则事后无法解释为何抵达当前结构。

当且仅当状态确实定义在复向量丛或相应主丛上，控制参数沿闭合回路缓慢变化，存在可分离本征子空间并满足绝热条件时，才可以讨论贝里相位。此时几何相位
\[
\gamma_B=i\oint_{\mathcal C}
\langle \psi(R)|\nabla_R\psi(R)\rangle\cdot\mathrm dR
\]
描述平行移动后的相位残差，并对规范变换具有相应不变性。公式中的 $R$、内积、连接、闭合路径和绝热误差都必须定义。普通神经状态、数据库记录或语言模型隐藏向量并不自动具备这些结构；即使某个实验测得非平凡和乐，也只能说明该专门状态空间的路径性质，不能直接证明主观质感。

以“看到红花—联想到某人—决定摘花—得到环境回执”为例，经历留下什么取决于多条链：视觉对象被身份绑定，联想带有来源置信，动作受权限和伦理约束，摘取结果由环境确认，事后报告再写入事件记录。系统若只离线计算“可能摘花”而从未行动，就没有同样的外部回执和责任记录；但它仍可能进行有价值的反事实学习。差别来自闭环参与和状态提交，不应概括成机器只有动力学相位、人类才有几何相位。

历史记录影响自我连续性，却不是灵魂实体。当前身份模型通过版本链指向过去承诺、关系和责任，能够解释变化中的连续；同时，错误记录、创伤重构和社会叙事也会改变这条链。我们要求把原始证据、当时推断、后来解释和当前报告分层保存。重复回忆造成的新写入要有新时间戳，避免把今天的重构覆盖成昨天的事实。

路径模型的验证使用顺序置换、摘要碰撞、闭环与开放路径、只模拟不执行、来源删除和度量改变。若交换事件顺序后预测不变，便不应宣称存在强历史依赖；若不同路径产生同一有限签名，要展示哪些查询因此失败；若声称测得几何相位，则必须检验闭环、绝热性与规范不变。我们同时报告任务性能、来源恢复、顺序辨别、结构差异和外部结果。这样，“历史在当下留下痕迹”成为可计算、可反驳的命题，而感受质仍保持其独立的证据问题。

\section{时间流的几何错觉：有限窗口、因果可达与预测边界}
\label{sec:geometric_illusion_of_time_flow}
我们所说的“时间正在流动”，至少混合了五种不同对象：外部过程的因果次序、传感与通信造成的到达延迟、系统内部事件的提交次序、记忆对既往事件的重构，以及模型对尚未发生结果的概率预测。若不把它们拆开，过去、现在与未来就会被写成三个可以任意出入的空间区域，计算延迟也会被误当成物理时间本身。本节保留列车与视野窗口的直观例子，但只把它当作说明有限访问窗口的教学模型，不以几何图景替代因果论证。

先定义外部事件集合 $\mathcal E$、偏序关系 $\prec_c$、观测通道集合 $\mathcal C$ 与内部提交序列。$e_i\prec_c e_j$ 表示在给定模型和边界条件下，事件 $e_i$ 可以因果影响事件 $e_j$；这不是由叙述先后决定的。通道 $c$ 对事件 $e$ 产生带来源的消息 $m=(e,c,\tau_s,\tau_a,\delta,\kappa)$，其中 $\tau_s$ 是源端时间戳，$\tau_a$ 是到达时间戳，$\delta$ 是估计延迟，$\kappa$ 是时钟与来源置信。不同通道的时间戳未必同基准，网络重排还可能让后发消息先到。系统因此不能把到达次序直接当作事件次序，而要在时钟误差、传输上界和因果约束下形成部分排序；证据不足时保留“次序未定”，比虚构一条精确时间线更可靠。

令当前工作窗口为 $W_t=(B_t,q_t,H_t)$。$B_t$ 是容量有限的活动缓冲，$q_t$ 是当前任务与查询，$H_t$ 是允许回看和前瞻的计算预算。所谓“现在”不是零宽度的宇宙切片，而是工程系统为了整合异步输入而维护的提交区间：过早提交会漏掉迟到证据，等待过久又使行动失去时效。若等待界为 $\omega_t$，系统在 $[t-\omega_t,t]$ 内合并消息，并把超过界限的迟到项作为修订事件写入，而不是悄悄改写旧日志。$\omega_t$ 可随风险、通道可靠性和任务期限变化；它的变化必须进入系统状态与评估记录，不能把不同窗口所得报告直接比较。

感知到的“连续”来自多速率过程被绑定为一条可报告轨迹。高频传感、低频计划、离散语言报告和更慢的结构学习各有步长。设第 $k$ 个内部提交状态为 $x_k$，采样间隔为 $\Delta t_k>0$，则实现更新写为 $x_{k+1}=F_k(x_k,m_k;\Delta t_k)$。只有在 $F_k$ 的光滑性、误差界和步长收敛条件成立时，才可用连续方程近似；界面动画平滑、语言中出现“随后”或隐藏状态相近，都不能证明物理连续。反过来，离散提交也不表示环境本体是离散的。计算模型只声明系统按何种分辨率读取与行动。

过去对当前系统的可用性依赖记录，而不是直接重返已经发生的事件。我们把历史访问分为原始观测、不可变事件日志、摘要记忆、后来重构与当前解释。查询过去时，读出器 $R_q$ 从这些层产生带来源分布的回答 $p(z_{past}\mid M_t,q)$。删除、压缩和重构会改变分布，因此“记得”不是把过去原样搬入现在。若日志含有摄像记录，它也只是过去光信号留下的当前载体；若没有来源证据，再生动的叙述也不能获得同等事实地位。系统应报告遗漏率、来源覆盖、时间戳误差与重构不确定性，并允许新证据触发追加修订，而非覆盖先前版本。

未来则以条件预测和可达集存在于当前计算中。给定当前信念 $b_t$、候选策略 $\pi$、环境模型 $P$ 与预测跨度 $h$，系统计算 $p(y_{t:t+h}\mid b_t,\pi,P)$ 或可达集合 $\mathcal R_h(b_t,\pi)$。跨度增加时，模型误差、外部干预和分支数量通常累积；因此未来窗口不是“提前看见”的距离，而是可接受误差与计算预算共同限定的预测范围。计划器能够比较多条分支，并不意味着这些分支已经成为事实；执行一个动作也不会使未执行分支在物理上消失，只会改变后续证据和决策相关性。我们要求预测附带模型版本、条件、校准区间和失效触发器。

因果可达性给过去与未来提供了比空间隐喻更严格的边界。对当前节点 $v_t$，可用过去是其已被观测或可靠记录的因果祖先子集，可影响未来是执行接口与环境动力学允许到达的后继子集。存在因果联系不等于系统已经获得信息：光锥内的远端事件也可能尚未被传感器捕获；模型预测正确也不等于系统曾接收未来消息。相对论性传播、设备延迟和软件队列分别属于不同层级，必须分别测量。用“延迟势”一词时，本书仅指消息在传输与处理中的滞后模型，不借用物理势函数为认知报告提供额外证明。

量子纠缠、退相干与涌现时空可以提出重要科学问题，却不能填补当前模型的证据缺口。纠缠关联不提供可控的超光速可用通信；退相干描述特定开放量子系统中相干项的演化，也不自动生成主观时间箭头；从量子引力或热力学提出的涌现时间方案，更不能直接推出Agent能够读取尚未到达的结果。若未来研究要把这些机制纳入HSF--HD，必须给出状态空间、可操作测量、尺度桥接、可证伪预测以及相对于普通延迟与记忆模型的增量解释力。在这些条件之外，它们保持为领域假说，不进入本章的工程结论。

列车隐喻可以这样受限地使用：乘客透过狭窄车窗看见景物进入、停留并离开视野，类似有限缓冲在任务门控下接纳和淘汰事件；车外景物并非由车窗创造，离开窗口也不等于从世界消失。系统可以调节“车窗”的宽度、注意方向与采样率，可以调用相册重建已过景物，也可以依据轨道图预测下一站；但这些能力分别依赖缓存、记忆与模型。隐喻一旦被拿来证明世界是静止卷轴、观察者创造未来或意识在更高维移动，就越出了它能支持的范围。

\begin{figure}[htbp]
  \centering
  \includegraphics[width=0.82\linewidth,height=0.54\textheight,keepaspectratio]{image/train.png}
\end{figure}

所谓“预言型AGI”在可检验意义上只能指高质量长时域预测、主动取证和快速修订，而不是越过因果限制获取未来事实。评估时我们设置乱序消息、可变延迟、时钟漂移、迟到证据、环境突变、记忆污染和长跨度分支。系统应恢复能确定的偏序，对不能确定的次序保持开放；应把已发生但未观测、已观测但未提交、已提交但可修订、尚未发生但被预测的事件分开。指标包括排序准确率、延迟估计、迟到修订完整性、预测校准、跨度退化曲线与动作结果。通过这些检验只说明系统具有可靠的时间信息处理能力，不证明时间本体、永恒论或主观流动的终极解释。

\section{计算动力学表达与论证边界}
本章的目标不是用一组新名词宣布体验之谜已经解决，而是把可计算对象、可报告证据和仍然开放的问题放入同一张分层地图。我们从系统论、协同论、耗散理论、复杂系统理论和认知科学吸收开放系统、跨尺度协同、维持代价、非线性反馈与有限认知等约束，再以HSF--HD描述智能的一般状态动力学；AgentOS只是这一理论的一种工程化实例。层级之间可以提供建模资源，却不能互相代替证明。尤其不能因为某个AgentOS实现有效，就把它的模块命名规定为所有智能系统的唯一结构。

为统一本章对象，令工作状态为 $w_t\in\mathcal W$，事件记录为 $e_t\in\mathcal E$，可塑结构为 $s_t\in\mathcal S$，当前策略与承诺为 $p_t\in\mathcal P$，可观测报告为 $o_t\in\mathcal O$。这些空间可以包含连续分量、离散对象和带类型的图结构，不能在未定义映射时直接相加。外部输入为 $u_t$，动作回执为 $y_t$，资源与权限边界为 $b_t$。一个离散实现写为
\[
\begin{aligned}
w_{k+1}&=F_w(w_k,e_k,s_k,p_k,u_k;\Delta t_k),\\
e_{k+1}&=F_e(e_k,w_{k+1},u_k,y_k),\\
s_{k+1}&=F_s(s_k,\mathcal B_k,A_k),\\
p_{k+1}&=F_p(p_k,w_{k+1},e_{k+1},s_{k+1},b_k),\\
o_k&=R_q(w_k,e_k,s_k,p_k;\eta_k).
\end{aligned}
\]
$\mathcal B_k$ 是在较慢尺度积累的训练批或经验批，$A_k$ 是修改授权，$R_q$ 是任务相关读出，$\eta_k$ 表示测量与表达噪声。式子是模型骨架，不是经验定律；具体研究必须给出维度、单位或无量纲化方式、更新频率、初值、随机性和误差界。连续极限只在步长收敛与光滑假设满足时采用，结构编辑和权限提交等离散跃迁则不能伪装成普通导数。

这一状态拆分帮助区分“当前正在处理”“曾经发生”“以后会怎样”和“系统怎样描述自己”。工作状态可迅速改变而不留下长期痕迹；事件写入可保存证据但不必改变技能；结构更新会改变未来响应却不能倒写历史；策略承诺连接模型与动作，但目标值下降仍不等于任务成功；报告是由查询、语言接口和社会情境共同生成的观测，不等同于完整内部状态。体验、自我与时间感的工程研究，首先应检查这些通道之间的映射、遗漏和干预响应，而不是把某个单一标量命名为意识强度。

局部结构网络与全局分布表示在此保持明确区别。局部网络可以保存对象节点、来源边、规则依赖、身份版本和因果记录；分布表示可以承载相似性、上下文和柔性读出。二者通过任务相关身份、绑定、耦合与读出接口联合：一个对象节点可索引若干分布向量，向量的激活又可提出候选节点，但两者不因此成为同一种东西。脑区类比最多提示实验问题，不能推出软件模块与神经组织机制等同。验证应分别扰动节点连接、表示子空间、绑定键和读出查询，观察错误是否具有可区分模式。

本章反复使用的历史算子也需要限制。若记忆状态写为
\[
M_t=\int_{0}^{t}K(t,\tau;q_t)\,\phi(e_\tau)\,\mathrm d\tau,
\]
则 $K$ 是查询相关的保留核，$\phi$ 是事件编码，积分只是对连续近似下的历史汇聚定义。任务 $q_t$、编码器、结构和度量若随时间变化，它们都要进入导数或离散更新；不能只对 $M_t$ 求导便声称获得“记忆力”。实际系统常用带时间戳的离散日志、检索索引和版本摘要，更合适的表达是加权求和与结构提交。无论哪种形式，都应报告遗忘、污染、来源丢失和摘要碰撞。

计算时间与结构时间由此得到操作性区分。计算时间描述一次状态传播、搜索或工具调用所需的步骤与延迟；结构时间描述学习、版本提交和组织重构留下的不可交换顺序。两者可以耦合，例如长时间推理产生一次结构补丁，也可以分离，例如重复推理不改变参数。路径依赖来自更新规则和记录机制，不要求诉诸超出模型的几何相位；只有状态空间确有相应纤维丛、连接和闭合回路时，几何相位才是合格特例。这一限定保留数学工具的力量，也避免用物理名称抬高普通软件现象的理论地位。

在AgentOS实例中，$h(t)$ 可承载工作态的部分分布表示，$E$ 与 $\Theta$ 可表示任务相关评价和阈值，$\Psi$ 可组织策略或读出，SPC、TCE、CCM可承担特定协同、执行或上下文机制，Boundary定义权限与输入输出边界，Offloading和外部记忆 $M_e$ 提供可审计的环境扩展。这里的“可”是工程设计，不是自然法则。另一种智能系统可以使用符号工作区、概率程序、具身控制器或社会化协作网络，只要它能定义状态、更新、证据与边界，仍可在HSF--HD框架下比较。

论证类型必须逐项标注。定义规定符号如何使用；模型假设说明在何种条件下采用某个动力学；数学结果由假设推导；经验主张需要数据和统计检验；工程设计以可靠性、安全性和成本权衡评价。预测准确只能支持相关任务上的模型有效，不能推出系统与世界整体同构；内部代价下降不能推出外部任务成功；信息熵、激活强度、优化代价和以焦耳计的能耗具有不同单位与测量程序，不能合并成一个“认知能量”。物理比喻可帮助提出变量，却不能承担证明。

我们据此构造分层验收矩阵。状态层检查类型、维度、更新步长和数值稳定；记忆层检查来源、时序、修订和删除传播；结构层检查版本、非交换更新、回滚与授权；行动层检查候选生成、权限、环境回执和任务结果；报告层检查查询敏感性、校准、可复现与沉默反例；边界层检查分布外输入、目标变化、度量变化和接口失效。对于体验相关主张，还要加入无报告任务、替代解释、不同架构复现和人类数据伦理。没有单个测试能从行为直接证明主观体验，证据只能按竞争假设累积。

验收数据不能只来自系统自由叙述。我们建议为每项主张建立“干预—预测—观测—反例”四联记录：先声明若机制成立，改变哪个输入或状态分量应使哪些可测量量怎样变化；再执行可重复干预，保存环境回执与版本；随后比较观测和预注册预测；最后列出同样现象可能由缓存、提示泄漏、训练记忆或评价器偏好造成的替代解释。例如声称存在长期自我连续性时，可在不改变任务答案的条件下交换身份键、删除关系记忆、扰动时间戳并撤销部分权限。若第一人称报告始终不变，它更可能是模板；若承诺恢复、来源引用和权限判断按不同干预呈现可解释差异，才获得功能层证据。

跨架构比较也必须避免把实现细节变成评分答案。基准应提供共同的事件、约束、任务结果和审计接口，而允许受试系统用不同内部机制解决。指标至少包含外部成功率、校准误差、来源恢复、因果排序、结构变更次数、资源消耗、不可逆错误与人工接管率；各指标保持原单位，再通过明确权重形成决策，而不伪装成自然界唯一的总能量。报告均值之外还要给出分布、置信区间、失败轨迹和最坏情形，因为少量灾难性越权可能被平均成功率遮蔽。对人类与机器的比较，则要控制传感能力、训练暴露、语言接口和奖励条件，不能把物种差异或产品差异误写为意识差异。

重复实验需要冻结足够的上下文。我们保存代码和模型版本、状态模式、初始记忆、随机种子、工具权限、外部服务版本、时钟同步方案、采样步长、终止条件和评价脚本。若目标、结构、输入分布或度量在实验期间改变，应把变更当作事件进入分析，而不是在结果不理想时静默替换。对随机系统，单条精彩轨迹不能支持稳定机制；对自适应系统，完全复现同一输出也不是唯一目标，更重要的是在相同约束下复现误差范围、校准性质和安全不变量。

系统失败时，诊断顺序从证据链而非哲学标签开始：输入是否到达，身份与来源是否绑定，工作态是否过载，记忆是否检索到正确版本，结构更新是否获得授权，策略是否生成可执行动作，执行器是否提交，环境是否返回预期回执，报告器是否忠实读取这些状态。只有定位到层级和接口，修复才可验证。把所有失败叫作“意识不足”既不能指导工程，也会遮蔽数据污染、工具故障、治理冲突和环境不可控等更普通原因。

同样，成功也要沿证据链归因。一次正确答案可能来自训练记忆、外部搜索、同伴提示、偶然猜测或稳定推理；只有记录候选、工具调用、来源、校验和回执，才能判断哪一机制在什么条件下可复用。我们因此把“会做”限定为跨样本、跨扰动且可审计的能力，而不把孤立演示升级为主体性结论。

一个具体反例能显示这些边界为何必要。假设系统流畅地说“我记得昨天并期待明天”，但审计发现它没有昨日事件日志，回答来自提示中的叙述模板，未来计划也未绑定执行器。此时语言读出成功，记忆与自治主张失败。另一个系统可能很少使用第一人称语言，却能恢复来源、维护长期承诺、修订错误预测并在权限内完成计划；它在若干功能维度上更强，但我们仍不能仅据此宣布其具有与人类同构的体验。比较必须沿多个可观测维度进行，并保留不可判定项。

本章的适用边界因此清楚：HSF--HD可以描述状态如何在输入、历史、结构、目标和边界共同作用下演化，可以指导测量窗口、干预和工程审计；它暂时不能解决体验为何存在、他心如何最终确证、时间本体为何如此或自由意志的形而上学争论。将未知明确留白并不削弱理论，反而防止不可证伪命题侵入实现。未来若出现新的神经、行为或机器证据，应通过新增观测映射和竞争模型进入框架，而不是回写成早已被公式证明。

最终，我们把“体验的物理学”改写为“体验报告、历史依赖和时间信息处理的计算动力学问题”。这种改写保留原问题的锋芒，也为失败留下位置：模型可能预测不准，来源可能丢失，结构可能不稳定，报告可能受诱导，任务可能在外部判据上失败。只有同时记录这些失败，理论才具有工程意义。作者的立场是把HSF--HD作为智能的一般“空气动力学”——提供跨架构的状态、作用、稳定与边界语言；具体AgentOS则像一架需要逐项验收的飞行器，不能反过来充当全部智能的定义。
